% V033 快记快懂快用 | 保留 summarybox 环境
\begin{summarybox}
\textbf{\textcolor{CSummary}{快记：一式、两层面积}}
\[
 \boxed{S_{\triangle OAB}=\tfrac12|x_1y_2-x_2y_1|},\qquad
 \boxed{S_{OADB}=|x_1y_2-x_2y_1|}.
\]
$A=(x_1,y_1),B=(x_2,y_2),O=(0,0),D=A+B$。
\textbf{为什么有 $\frac12$？} $AB$ 是平行四边形 $OADB$ 的对角线，等分面积。

\medskip
\textbf{\textcolor{CSummary}{快懂：四种辅助线，任选起点与方向}}
\begin{itemize}[leftmargin=*,itemsep=0.25em]
\item 过 $B$ 竖直交 $OA$：$y_{E_v}=x_2y_1/x_1$，$S=\frac12|x_1|BE_v$。
\item 过 $B$ 水平交 $OA$：$x_{E_h}=x_1y_2/y_1$，$S=\frac12|y_1|BE_h$。
\item 过 $A$ 竖直交 $OB$：$y_{F_v}=x_1y_2/x_2$，$S=\frac12|x_2|AF_v$。
\item 过 $A$ 水平交 $OB$：$x_{F_h}=x_2y_1/y_2$，$S=\frac12|y_2|AF_h$。
\end{itemize}
\textbf{统一检查：} 相应分母必须非零；交点在对边内部时面积相加，在延长线上时相减。
若分母为零，直接以水平或竖直的 $OA$ 或 $OB$ 为底；若对应边退化为点，面积为零。

\medskip
\textbf{\textcolor{CSummary}{快用：考场识别信号与方法选择}}
\begin{itemize}[leftmargin=*,itemsep=0.25em]
 \item \textbf{求得最快：} 已知两向量坐标 $\Rightarrow$ 直接代入二阶行列式面积。
 \item \textbf{需要讲清来源：} 优先用竖直/水平共底面积法，初中知识即可理解。
 \item \textbf{标准解析法：} 选 $OA$ 为底，求 $B$ 到 $OA$ 的距离；分母相消。
 \item \textbf{任意三点：} 先平移到原点，再代面积公式。
 \item \textbf{三点共线：} 两倍面积 $x_1y_2-x_2y_1=0$；无需事先判定。
 \item \textbf{关联 C088：} $y_1=(b/a)x_1$、$y_2=-(b/a)x_2$
 $\Rightarrow\ S=\frac ba|x_1x_2|$。
\end{itemize}

\medskip
\textbf{\textcolor{CSummary}{高价值变式}}
若 $A\in\{y=m_1x\}$、$B\in\{y=m_2x\}$，
\[
 \boxed{S_{\triangle OAB}=\frac12|m_1-m_2|\,|x_1x_2|.}
\]
\textbf{考前检查：} 交叉相乘、绝对值、三角形的 $\frac12$、延长线面积相减、坐标平移。
\end{summarybox}
